% Section 19: Open Problems
\section{Open Problems and Future Directions}
\label{sec:open_problems}

\subsection{Mathematical Challenges}

Despite the progress documented in this work, several significant mathematical challenges remain:

\subsubsection{Type III Factor Structure}

\begin{problem}
Prove that local algebras $\mathcal{A}(\mathcal{O})$ in CTUFT are type III$_1$ factors and establish the split property.
\end{problem}

\textit{Status}: The type III$_1$ property is expected based on general quantum field theory results, but a rigorous proof for the torsion field requires showing the modular spectrum is $\mathbb{R}_+$.

\subsubsection{Non-Perturbative Construction}

\begin{problem}
Construct the interacting Wightman functions non-perturbatively, potentially using lattice regularization or constructive field theory methods.
\end{problem}

\textit{Status}: Current results are perturbative. A lattice formulation exists but the continuum limit requires further analysis.

\subsubsection{Infinite-Dimensional Measure Theory}

\begin{problem}
Develop rigorous measure theory for the infinite-dimensional phase space in geometric quantization.
\end{problem}

\textit{Status}: Finite-mode truncations are well-defined, but the limit requires careful treatment of the infinite-dimensional Wiener measure.

\subsection{Physical Extensions}

\subsubsection{Quantum Information Theory}

\begin{problem}
Compute entanglement entropy using the replica trick in the algebraic framework and relate to black hole entropy.
\end{problem}

\subsubsection{Cosmological Applications}

\begin{problem}
Extend the rigorous framework to curved spacetime backgrounds, particularly inflationary geometries.
\end{problem}

\subsection{Interdisciplinary Connections}

\subsubsection{Noncommutative Geometry}

\begin{problem}
Establish the precise relationship between CTUFT and Connes' noncommutative geometry program.
\end{problem}

\subsubsection{String Theory}

\begin{problem}
Investigate whether CTUFT emerges as a low-energy limit of string theory or represents a distinct UV completion.
\end{problem}

\subsection{Experimental Implications}

The resolution of these mathematical problems has direct physical consequences:
\begin{itemize}
    \item Type III structure $\to$ entanglement properties of the vacuum
    \item Non-perturbative effects $\to$ corrections to scattering amplitudes
    \item Curved spacetime extension $\to$ cosmological predictions
\end{itemize}
